\section{Fast-Charging Optimization Problem}
\label{sec:problem}

\subsection{Protocol Parameterization}

We optimize a three-stage constant-current charge from
$s_0=0$ to $s_{\mathrm{tar}}=0.80$. The decision vector is
\begin{equation}
\btheta=[I_1,I_2,I_3,\Delta s_1,\Delta s_2]^\top\in\Theta ,
\label{eq:theta}
\end{equation}
where $I_k$ is a current coordinate defined for a 5-Ah reference cell and
$\Delta s_1,\Delta s_2$ are the first two SOC increments.
Table~\ref{tab:variables} gives the box $\Theta$. For a parameterization with
effective capacity $Q_{\mathrm{eff}}$, the applied current is
$I_k^{\mathrm{app}}=I_kQ_{\mathrm{eff}}/(5~\mathrm{Ah})$. Thus, the
implementation preserves the reference C-rate when the PyBaMM
parameterization changes.

\begin{table}[!t]
\caption{Three-Stage Protocol Search Space}
\label{tab:variables}
\centering
\begin{tabular}{lccc}
\toprule
Variable & Symbol & Bounds & Coordinate \\
\midrule
Stage-1 current & $I_1$ & $[2.0,\;6.0]$ & A (5-Ah ref.) \\
Stage-2 current & $I_2$ & $[2.0,\;5.0]$ & A (5-Ah ref.) \\
Stage-3 current & $I_3$ & $[2.0,\;3.0]$ & A (5-Ah ref.) \\
Stage-1 SOC increment & $\Delta s_1$ & $[0.10,\;0.40]$ & -- \\
Stage-2 SOC increment & $\Delta s_2$ & $[0.10,\;0.30]$ & -- \\
\bottomrule
\end{tabular}
\end{table}

The final increment and ideal stage times are computed together as
\begin{equation}
\Delta s_3=0.80-\Delta s_1-\Delta s_2,\qquad
t_c=\sum_{k=1}^{3}\frac{3600Q_{\mathrm{eff}}\Delta s_k}
{I_k^{\mathrm{app}}}.
\label{eq:stages_time}
\end{equation}
The hard condition $\Delta s_1+\Delta s_2\leq0.70$ retains at least
0.10 SOC for the final stage. The selector also uses
$\Delta s_1+\Delta s_2\leq0.65$ and $I_1\geq I_2\geq I_3$ as soft
initialization preferences, not as physical or feasibility guarantees.

\subsection{Objectives and Encoded Constraints}

For a successfully simulated protocol, the constrained minimization problem is
\begin{equation}
\begin{aligned}
\min_{\btheta\in\Theta}\quad&
\mathbf f(\btheta)=
[t_c(\btheta),\Delta T_p(\btheta),D_{\mathrm{chg}}(\btheta)]^\top\\
\mathrm{s.t.}\quad&
\Delta s_1+\Delta s_2\leq0.70,\quad
V_{\mathrm{peak}}(\btheta)\leq V_{\max},\\
&T_{\mathrm{peak}}(\btheta)\leq T_{\max}.
\end{aligned}
\label{eq:cmop}
\end{equation}
Here $\Delta T_p$ is the scaled peak temperature rise and
$D_{\mathrm{chg}}$ is the protocol-level degradation proxy in arbitrary units
(a.u.). Solver failures or encoded-limit violations receive the archived
penalty vector. Chen2020 uses $V_{\max}=4.4$~V and
$T_{\max}=328.15$~K; Ecker2015 uses 4.2~V and 333.15~K. These values define the
reproduced simulations and are not general charging recommendations.

\subsection{Electrochemical--Thermal Evaluation Model}
\label{sec:model}

The evaluator uses the single-particle model with electrolyte (SPMe)
~\cite{ref11} implemented in PyBaMM~\cite{ref12}. We use the Chen2020
parameterization for the 5-Ah LG INR21700-M50 cell~\cite{ref13} and the
PyBaMM Ecker2015 parameterization~\cite{ref14}. The latter has an effective
nominal capacity of 0.15625~Ah in this software configuration; this value is
not the capacity of the cell used in the original identification experiment.

The lumped thermal state satisfies~\cite{ref22}
\begin{equation}
\begin{aligned}
m c_p\frac{\mathrm dT}{\mathrm dt}
&=\dot Q_{\mathrm{gen}}-hA_s(T-T_{\mathrm{amb}}),\\
\Delta T_p&=1.10\,[\max_tT(t)-298.15~\mathrm{K}].
\end{aligned}
\label{eq:thermal}
\end{equation}
The factor 1.10 is applied identically to all archived methods. It is an
implementation alignment factor rather than a cell-specific calibration.

\subsection{Protocol-Level Degradation Proxy}

The archived empirical calculation can be written explicitly as
\begin{align}
\bar T_a&=298.15+0.61(\bar T-298.15),\nonumber\\
\tau_{\mathrm{emp}}
&=(2896.6\bar s_{\%}+7411.2)
\exp\!\left(\frac{-31500+152.5\bar I}{8.314\bar T_a}\right),\nonumber\\
C_{\mathrm{emp}}&=\left(\frac{20}{\tau_{\mathrm{emp}}}\right)^{1/0.57},
\label{eq:capacity_proxy}\\
D_{\mathrm{chg}}
&=100\,\frac{Q_{\mathrm{eff}}}{C_{\mathrm{emp}}}.
\label{eq:aging}
\end{align}
The trajectory averages use SOC in percent, temperature in kelvin, and current
in amperes. The form is inherited from a control-oriented empirical
cycle-life relation~\cite{ref24}, while the temperature factor 0.61 and all
coefficients remain uncalibrated for Chen2020 and Ecker2015. Consequently,
$D_{\mathrm{chg}}$ is used only to rank protocols generated by the same
evaluator. It is neither a measured percentage nor a calibrated prediction of
cycle life.
